\section{Introduction}
\label{sec:intro}


In daily life, humans are capable of simultaneously perceiving the visual and auditory information around them. After processing this information through the brain, they express feedback through writing, vocalization, or using tools (and physical actions), thereby engaging in information exchange with various organisms in the world and exhibiting intelligence. In recent years, general artificial intelligence has become increasingly visible, largely due to advancements in Large Language Models (LLMs)~\citep{gpt3,gpt4,gpt4o,gemini,claude,claude2,claude3,qwen,qwen2,llama,llama2,llama3}. These models, trained on vast amounts of textual data, represent high-level discrete representation created by humans, showcasing the ability to solve complex problems and learn rapidly. Furthermore, in the realm of understanding, Language-Audio-Language Models (LALMs)~\citep{gpt4o,anonymous2023salmonn,qwenaudio, qwen2-audio} and Language-Visual-Language Models (LVLMs)~\citep{blip2,llava,instructblip,minigpt-4,kosmos,qwenvl,llava1.5,cogvlm,gpt4v,gemini} have helped LLMs to further extend auditory and visual capabilities in an end-to-end manner. %Recently, there has been rapid development in end-to-end speech generation, and intelligent speech dialogue capabilities have demonstrated enormous potential. 
However, efficiently unifying all these different understanding modalities in an end-to-end fashion, utilizing as much data as possible, and providing responses in both text and speech streams akin to human communication still presents a significant challenge.


The development of a unified and intelligent omni-model requires careful consideration of several key factors. First, it is crucial to implement a systematic method for the joint training of various modalities, including text, images, videos, and audio, to foster mutual enhancement among them. This alignment is particularly important for video content, where synchronization of the temporal aspects of audio and visual signals is necessary. Second, it is essential to manage potential interference among outputs from different modalities, ensuring that the training processes for outputs such as text and voice tokens do not disrupt each other. Finally, there is a need to explore architectural designs that enable real-time understanding of multimodal information and allow for efficient audio output streaming, thereby reducing initial latency.


In this report, we introduce \method, a unified single model capable of processing multiple modalities and generating text and natural speech responses simultaneously in a streaming format. To tackle the first challenge, we propose a novel position embedding approach, named \textbf{TMRoPE}~(\textbf{T}ime-aligned \textbf{M}ultimodal \textbf{RoPE}). We organize these audio and video frames in an interleaved structure to represent video sequences in time order. 
For the second challenge, we present Thinker-Talker architecture, wherein Thinker is tasked with text generation while the Talker focuses on generating streaming speech tokens. Talker receives high-level representations directly from Thinker. This design is inspired by the way humans utilize different organs to produce various signals, which are simultaneously coordinated through the same neural networks. As a result, Thinker-Talker architecture is end-to-end jointly trained, with each component dedicated to generating distinct signals.
To address the challenges associated with streaming and to facilitate the pre-filling necessary for real-time comprehension of multimodal signals, we propose modifications to all multimodal encoders by adopting a block-wise streaming processing approach. In order to support streaming speech generation, we implement a dual-track autoregressive model that generates speech tokens, alongside a DiT model which converts these tokens into waveforms, thereby enabling streaming audio generation and minimizing initial latency. This design aims to enable the model to process multimodal information in real-time and effectively perform pre-filling, thereby enabling the concurrent generation of text and speech signals. 

\method is comparable with the similarly sized Qwen2.5-VL~\citep{qwen2vl} and outperforms Qwen2-Audio~\citep{qwen2-audio} in image and audio capabilities respectively. Furthermore, \method achieves state-of-the-art performance on multimodal benchmarks such as OmniBench~\citep{li2024omnibench} and AV-Odyssey Bench~\citep{gong2024av}. Notably, \method's performance in end-to-end speech instruction following is comparable to its capabilities with text inputs, as evidenced by benchmarks such as MMLU~\citep{mmlu} and GSM8K~\citep{gsm8k}. As for speech generation, \method achieves 1.42\%, 2.33\% and 6.54\% WER on seed-tts-eval~\citep{seedtts} test-zh, test-en and test-hard set respectively, outperforming MaskGCT~\citep{maskgct} and CosyVoice~2~\citep{cosyvoice2}.

The key features of \method can be summarized as:

\begin{itemize}
    \item  We introduce \method, a unified model that can perceive all modalities and simultaneously generate text and natural speech responses in a streaming fashion.
    \item We present a novel positional embedding algorithm, termed TMRoPE, which explicitly incorporates temporal information for synchronizing audio and video.
    \item We propose the Thinker-Talker Architecture to facilitate real-time comprehension and speech generation.
    \item \method demonstrates strong performance across all modalities when benchmarked against similarly sized single-modality models. It significantly enhances the capability of following voice commands, achieving performance levels comparable to pure text input. For tasks that involve integrating multiple modalities, such as those evaluated in OmniBench~\citep{li2024omnibench}, \method achieves state-of-the-art performance. Notably, \method achieves strong performance on seed-tts-eval~\citep{seedtts}, demonstrating robust speech generation abilities.
\end{itemize}
